\subsection{Nonunique decompositions of projectives: Problem 4.55}
\label{cand:4.55}
\problemstatement{4.55}

Let \(G=3.A_7\), the perfect central cover of order \(7560\),
and \(R=\Z_{(5)}\). There is an \(RG\)-projective module of
\(R\)-rank \(55{,}440\) with two indecomposable decompositions
having no isomorphism class of summand in common. This addresses
uniqueness itself, beyond the earlier failure of the condition
that indecomposable endomorphism rings be local.

We first give the descent facts used. If an \(R\)-free
\(G\)-lattice \(L\) has projective reduction modulo five,
its restriction to a Sylow \(5\)-subgroup \(S\) is free:
lift a free basis over the local algebra \(\F_5S\),
apply Nakayama's lemma to the resulting \(RS\)-map, and
use equal \(R\)-ranks for injectivity. Averaging over a
transversal, with \([G:S]\) invertible in \(R\), splits
\(RG\otimes_{RS}L\to L\). Thus \(L\) is projective.
Also, two finite projective \(RG\)-modules with isomorphic
reductions are isomorphic: projectivity lifts the reduction
map, Nakayama gives surjectivity, and equal \(R\)-ranks
give injectivity.

Suppose a finite projective \(\Z_5G\)-module \(M\) has
generic fiber obtained from a rational \(G\)-module \(V\).
In their identified \(\Q_5\)-space, approximate a
\(\Z_5\)-basis of \(M\) by rational vectors modulo \(5M\).
They form both a \(\Z_5\)-basis of \(M\) and a
\(\Q\)-basis of \(V\). The \(G\)-matrices in this basis
have entries in \(\Q\cap\Z_5=R\). Their \(R\)-span
is a stable lattice completing to \(M\), and is
projective by the preceding reduction criterion.

The finite character input is the following part of the
\(5\)-modular decomposition data for \(3.A_7\)
\cite{a7-decomposition,ctbl-3a7}. Use four Frobenius pairs
of absolute Brauer characters, denoted \(A,B,C,D\),
of respective absolute degrees \(3,6,18,21\).
Their \(\F_5\)-simple modules have dimensions
\(6,12,36,42\), and their projective covers have
dimensions \(90,60,90,90\). For one faithful central
character, the relevant rows are
\[
\begin{array}{c|r|rrrr}
 \text{GAP ordinary index}&\text{degree}&A&B&C&D\\ \hline
 10&6 &0&1&0&0\\
 16&21&0&0&0&1\\
 18&21&1&0&1&0\\
 20&24&1&0&0&1\\
 22&24&0&1&1&0
\end{array}
\]
and the conjugate rows have indices \(11,17,19,23,21\).
These indices use \artifact{CharacterTable("3.A7")} and its 5-modular
table in GAP~4.16.1 with the character-table library \texttt{CTblLib}
version~1.3.11. All other rows are zero in these columns.
The ordinary Galois orbits are
\(\{10,11\},\{16,17\},\{18,19\},\{20,21,22,23\}\).
Indices refer to the retained GAP table, not a numbering
across all covers in the printed tables.

Brauer reciprocity shows that a completed projective
whose reduction has cover multiplicities \((a,b,c,d)\)
has generic ordinary multiplicities \(b,d,a+c\) on
the first three paired orbits, and \(a+d,b+c\) on
the two halves of the degree-24 orbit. Its character
is rational-valued exactly when
\begin{equation}\label{eq:four-rays}
 a+d=b+c.
\end{equation}
This is necessary for every rationally defined
projective with this support.

Consider the four rays
\[
 r_1=(1,1,0,0),\quad r_2=(1,0,1,0),\quad
 r_3=(0,1,0,1),\quad r_4=(0,0,1,1).
\]
Each gives a completed projective, since idempotents
lift over the complete ring \(\Z_5G\), and each
satisfies \eqref{eq:four-rays}. Multiplication of
its generic character by \(N=168\) gives a rational
form without assuming Schur index one. Indeed, for
an ordinary Galois orbit \(\mathcal O\), the sum
of its primitive complex central idempotents lies
in \(\Q G\), and \(\Q Ge_{\mathcal O}\) has character
\(\sum_{\chi\in\mathcal O}\chi(1)\chi\).
Since \(168\) is divisible by \(6,21,24\), the
desired multiple is a nonnegative sum of these
rational characters. Semisimple \(\Q_5G\)-modules
are determined by their characters after scalar
extension, so the descent argument gives projectives
\(L_i\) over \(RG\) with reduction vectors \(168r_i\).

Now \(r_1+r_4=r_2+r_3\), and reduction detects
isomorphism of projectives. Hence
\[
 L_1\oplus L_4\cong L_2\oplus L_3.
\]
The common rank is \(168(150+180)=55{,}440\).
Every nonzero direct summand of \(L_1\) has reduction
supported on \(A,B\); finite-dimensional
Krull--Schmidt and \eqref{eq:four-rays} force its
vector onto the positive ray through \(r_1\).
The same holds for each of the other three rays.
They meet only at zero. Thus no nonzero summand
on one side can be isomorphic to a summand on
the other. Splitting successively terminates by
positive finite \(R\)-rank, giving incompatible
indecomposable refinements.

The character data are a genuine dependency: the
four-ray calculation alone would not construct
modules. The archive reconstructs the ordinary
table from the actual group and the faithful
decomposition rows from modular representations,
in addition to comparison with published tables.
The \(L_i\) themselves are not asserted to be
indecomposable. Source:
\artifact{research/4.55-proof.md}.

